% Раздел предназначен для подключения через \input.
% Требуются стандартные пакеты amsmath и amssymb,
% а также обозначения проекта: \betav, \fcl, \Imav, \UV, \JJJ, \dimd.

\section{Variable projection и Riemannian L-BFGS для single-index задачи}
\label{sec:varpro-rlbfgs}

На фиксированном внешнем шаге структурной адаптации считаются заданными
матрицы $\UV_j\in\mathbb R^{n_{\Phi,j}\times\dimd}$ и векторы
$\Imav_j\in\mathbb R^{n_{\Phi,j}}$. Индекс внешнего шага далее опускается.
Рассматривается стабилизированная задача
\begin{equation}
\min_{\substack{\|\betav\|=1\\(\fcl_j)_{j\in\JJJ}}}
\frac12\sum_{j\in\JJJ}
\left(
 \|\Imav_j-\fcl_j\UV_j\betav\|^2
 +\eta_j\fcl_j^2
\right),
\qquad \eta_j\ge 0.
\label{eq:rlbfgs-original}
\end{equation}
При $\eta_j=0$ получается исходная задача. Небольшое $\eta_j>0$ рекомендуется только для локальных блоков, в которых $\|\UV_j\betav\|$ может быть близка к нулю.

\subsection{Исключение локальных коэффициентов}

Для фиксированного $\betav$ введём
\begin{equation}
 \boldsymbol{q}_j=\UV_j\betav,
 \qquad
 a_j=\langle\Imav_j,\boldsymbol{q}_j\rangle,
 \qquad
 b_j=\|\boldsymbol{q}_j\|^2+\eta_j.
\end{equation}
Точный минимум по локальному коэффициенту равен
\begin{equation}
 \widehat{\fcl}_j(\betav)=\frac{a_j}{b_j}.
\label{eq:profiled-slope}
\end{equation}
После подстановки \eqref{eq:profiled-slope} остаётся задача на единичной сфере
$\mathbb S^{\dimd-1}$:
\begin{equation}
 \min_{\|\betav\|=1}F(\betav),
 \qquad
 F(\betav)=
 \frac12\sum_{j\in\JJJ}
 \left(
  \|\Imav_j\|^2-\frac{a_j^2}{b_j}
 \right).
\label{eq:profiled-objective}
\end{equation}
Таким образом, Riemannian L-BFGS оптимизирует только $\dimd$-мерный вектор
$\betav$, а все коэффициенты $\fcl_j$ пересчитываются аналитически.

\subsection{Градиент на сфере}

По теореме об огибающей евклидов градиент профилированной функции можно
вычислять без дифференцирования формулы
$\widehat{\fcl}_j(\betav)$:
\begin{equation}
 \nabla F(\betav)
 =-
 \sum_{j\in\JJJ}
 \widehat{\fcl}_j\,
 \UV_j^{\mathsf T}
 \left(
  \Imav_j-\widehat{\fcl}_j\UV_j\betav
 \right).
\label{eq:euclidean-gradient}
\end{equation}
Касательное пространство сферы в точке $\betav$ имеет вид
\begin{equation}
 T_{\betav}\mathbb S^{\dimd-1}
 =\{\boldsymbol{\xi}\in\mathbb R^{\dimd}:\langle\betav,\boldsymbol{\xi}\rangle=0\}.
\end{equation}
Поэтому риманов градиент равен ортогональной проекции
\begin{equation}
 \operatorname{grad}F(\betav)
 =\nabla F(\betav)
 -\betav\langle\betav,\nabla F(\betav)\rangle.
\label{eq:riemannian-gradient}
\end{equation}
Формулы \eqref{eq:profiled-objective}--\eqref{eq:riemannian-gradient}
требуют только умножений $\UV_j\betav$ и
$\UV_j^{\mathsf T}\boldsymbol{r}_j$. Матрицы $\UV_j^{\mathsf T}\UV_j$ и матрицы
размера $\dimd\times\dimd$ не строятся.

\subsection{Шаг Riemannian L-BFGS}

Пусть $\boldsymbol{g}_t=\operatorname{grad}F(\betav_t)$. По последним $m$ парам
$(\boldsymbol{s}_i,\boldsymbol{y}_i)$ двухцикловая рекурсия L-BFGS строит приближение действия
обратного риманова гессиана и направление
\begin{equation}
 \boldsymbol{d}_t=-\mathcal H_t\boldsymbol{g}_t
 \in T_{\betav_t}\mathbb S^{\dimd-1}.
\end{equation}
Если $\langle\boldsymbol{g}_t,\boldsymbol{d}_t\rangle\ge0$, используется безопасный шаг
$\boldsymbol{d}_t=-\boldsymbol{g}_t$.

Длина шага $\alpha_t$ выбирается обратным поиском из условия Армихо
\begin{equation}
 F\bigl(R_{\betav_t}(\alpha_t\boldsymbol{d}_t)\bigr)
 \le
 F(\betav_t)
 +c_1\alpha_t\langle\boldsymbol{g}_t,\boldsymbol{d}_t\rangle.
\label{eq:armijo}
\end{equation}
Для сферы применяется нормирующая ретракция
\begin{equation}
 R_{\betav}(\boldsymbol{\xi})
 =\frac{\betav+\boldsymbol{\xi}}{\|\betav+\boldsymbol{\xi}\|}.
\label{eq:sphere-retraction}
\end{equation}
После получения $\betav_{t+1}$ касательные векторы переносятся простой
проекцией
\begin{equation}
 \mathcal T_{t\to t+1}(\boldsymbol{z})
 =\boldsymbol{z}-\betav_{t+1}
   \langle\betav_{t+1},\boldsymbol{z}\rangle.
\label{eq:vector-transport}
\end{equation}
Пары для обновления памяти определяются как
\begin{equation}
 \boldsymbol{s}_t=\mathcal T_{t\to t+1}(\alpha_t\boldsymbol{d}_t),
 \qquad
 \boldsymbol{y}_t=\boldsymbol{g}_{t+1}-\mathcal T_{t\to t+1}(\boldsymbol{g}_t).
\label{eq:curvature-pairs}
\end{equation}
Пара сохраняется только при выполнении условия положительной кривизны
\begin{equation}
 \langle\boldsymbol{s}_t,\boldsymbol{y}_t\rangle
 >\varepsilon_{\mathrm{curv}}\|\boldsymbol{s}_t\|\,\|\boldsymbol{y}_t\|.
\label{eq:curvature-test}
\end{equation}
Это предотвращает неустойчивое обновление аппроксимации гессиана.

\subsection{Вычислительная схема}

Одна итерация состоит из следующих операций:
\begin{enumerate}
 \item для всех $j$ вычислить $\boldsymbol{q}_j=\UV_j\betav_t$,
       $\widehat{\fcl}_j$, значение $F(\betav_t)$ и остатки;
 \item вычислить евклидов и риманов градиенты по
       \eqref{eq:euclidean-gradient} и \eqref{eq:riemannian-gradient};
 \item получить направление двухцикловой рекурсией L-BFGS;
 \item выполнить поиск шага и ретракцию \eqref{eq:sphere-retraction};
 \item перенести касательные векторы и обновить ограниченную память.
\end{enumerate}

Если $p_j$ обозначает число строк $\UV_j$, стоимость вычисления функции и
градиента составляет
\begin{equation}
 O\!\left(\sum_{j\in\JJJ}p_j\dimd\right).
\end{equation}
Дополнительная память оптимизатора равна $O(m\dimd)$, не считая хранения
$\Imav_j$ и $\UV_j$. Сумма по центрам может вычисляться батчами или
потоково, поэтому алгоритм не требует хранения дополнительных тензоров,
зависящих от числа центров.

Остановка выполняется по норме риманова градиента
$\|\operatorname{grad}F(\betav_t)\|$ и относительному изменению
профилированной функции. После остановки локальные коэффициенты
восстанавливаются по \eqref{eq:profiled-slope}. Из-за эквивалентности
$(\betav,\fcl_j)$ и $(-\betav,-\fcl_j)$ знак итогового направления можно
согласовать с начальной оценкой.
